\chapter{What the sandbox takes away}\label{ch:sandbox}

ATE removes the Lua entry points that could run code it has not checksummed:
\texttt{dofile}, \texttt{loadfile} and \texttt{load} are gone, \texttt{package.path} and
\texttt{package.cpath} are empty, \texttt{debug} is reduced to \texttt{debug.traceback}, and
there is no \texttt{python} table.

This is traceability, not security. \texttt{io} and \texttt{os} are still there --- a schedule
can read a file if it knows the path. What ATE is guaranteeing is that it knows exactly which
files a Test Run ran, so a report can name them and their checksums. A schedule written in
Lua, or one written in a hurry, is trusted; a schedule that reaches for code outside its Media
Release is not, because the Test Definition cannot describe it.

\noindent\textbf{\texttt{sandbox.tslua}}\par
\lstinputlisting[language=ATELua]{../examples/test-writer/sandbox.tslua}

\section{Where each call is legal}

The top level of a \texttt{.tslua} file runs before a Test Run exists. Only the builders,
\texttt{require}, \texttt{ATE.log} and the enums are legal there. Everything else ---
\texttt{ATE.open}, \texttt{ATE.report}, \texttt{ATE.input}, \texttt{ATE.addEquipment} --- raises
a load error that names the function:

\begingroup\footnotesize
\begin{verbatim}
forbidden.tslua:7: ATE.open: only legal after the schedule has loaded
\end{verbatim}
\endgroup

That is why the Device is opened in a Pre-test Action and not at the top of the file.

\noindent\textbf{\texttt{forbidden.tslua}}\par
\lstinputlisting[language=ATELua]{../examples/test-writer/forbidden.tslua}

In a Pre-test or Post-test Action everything is legal except \texttt{ATE.input} and
\texttt{ATE.report}, which belong to steps. In a step action, everything is legal.

\cando{tell which Lua and \texttt{ATE} calls a schedule may make, and where, and read the
error a schedule gets when it makes one in the wrong place.}
